% Generated by Listas/make_statements.py from pratica-lecture-01.md, item 24. Do not edit.
\textbf{Source.} Tadelis (2013), \emph{Game Theory: An Introduction}, Exercise 4.4 ``eBay's Recommendation'', pp.~76--77 {[}PDF 95--96{]}. Reproduced verbatim. Answer in the source: none in the book (the preface, p.~xiii, says about half are solved on Princeton UP's website; not checked).



\begin{quote}

\textbf{eBay's Recommendation:} It is hard to imagine that anyone is not familiar with eBay, the most popular auction web site by far. In a typical eBay auction a good is placed for sale, and each bidder places a ``proxy bid,'' which eBay keeps in memory. If you enter a proxy bid that is lower than the current highest bid, then your bid is ignored. If, however, it is higher, then the current bid increases up to one increment (say, \$0.01) above the \emph{second highest} proxy bid. For example, imagine that three people have placed bids on a used laptop of \$55, \$98, and \$112. The current price will be at \$98.01, and if the auction ended the player who bid \$112 would win at a price of \$98.01. If you were to place a bid of \$103.45 then the player who bid \$112 would still win, but at a price of \$103.46, while if your bid was \$123.12 then you would win at a price of \$112.01.



Now consider eBay's historical recommendation that you think hard about the value you impute to the good and that you enter your true value as your bid---no more, no less. Assume that the value of the good for each potential bidder is independent of how much other bidders value it.



\begin{enumerate}

\def\labelenumi{\alph{enumi}.}

\item

  Argue that bidding more than your valuation is weakly dominated by actually bidding your valuation.

\item

  Argue that bidding less than your valuation is weakly dominated by actually bidding your valuation.

\item

  Use your analysis to make sense of eBay's recommendation. Would you follow it?

\end{enumerate}

\end{quote}



\textbf{Optional extras.}

- Muñoz-Garcia Ch. 1 Ex. 6--8 (pp.~11--14) and Ex. 11 (p.~20, three players): more IESDS drill, all with answers.

- Muñoz-Garcia Ch. 1 Ex. 17 (p.~33): a second weak-deletion order example, after item 19.

- Tadelis 3.1, 3.3, 3.5 (p.~58): the other verbal-to-matrix stories, if items 1--2 felt slow.



\begin{center}\rule{0.5\linewidth}{0.5pt}\end{center}

